\section[A Hermitian symmetric space in signature (4n,4n) with compact conformal quotients]{A Hermitian symmetric space in signature $\boldsymbol{(4n,4n)}$ \\ with compact conformal quotients}
 In the following, we use the index conventions
\[A,B,C,\ldots \in \{1,\ldots, 4n\},\qquad a,b,c,\ldots\in \{1\ldots ,2n\},\qquad i,j,k,\ell,\ldots\in \{1,\ldots, n\}.\] We use the Einstein summation convention for these ranges of indices.

Let $(\sigma_{ij})$ be a symmetric complex $n\times n$ matrix and define the quadratic polynomial
\[\sigma\bigl(z^1,\ldots, z^{2n}\bigr):= \sigma_{ij}z^{i}z^{j+n},\]
on $\C^{2n}$.
Using $\sigma$, we define
 $f \colon \C^{4n}\to \rr$ as
\begin{align}
f\bigl(z^1,\ldots , z^{4n}\bigr)
&{}=
\delta_{ab} z^a\ol{z}^{b+2n}+
\sigma\bigl(z^1,\ldots, z^{2n}\bigr)\overline{\sigma\bigl(z^1,\ldots, z^{2n}\bigr)}
\nonumber\\
&{}=
\delta_{ab} z^a\ol{z}^{b+2n}
+
\sigma_{ij}\bar{\sigma}_{k\ell}\bigl(
 z^{i}z^{j+n}\ol{z}^{k} \ol{z}^{\ell+n}\bigr).\label{potential}
\end{align}
As K\"{a}hler potential, this function defines an indefinite K\"ahler metric of signature $(4n,4n)$ on $\C^{4n}=\rr^{8n}$ by
\begin{gather}\label{metric}
g= h_{AB} \,
 \d z^A \cdot \d \ol{z}^B,\qquad \text{where} \quad h_{AB}=\frac{\partial^2 f}{\partial z^A\partial \ol{z}^B},
 \end{gather}
with K\"{a}hler form
 \begin{align*}
\omega
=
\frac{\mathrm{i}}{2} h_{AB}\,
 \d z^A \wedge \d \ol{z}^B.
\end{align*}
The only non-vanishing terms in $h_{AB}$ are the following:
\begin{gather}
h_{a\, b+2n} = h_{a+2n\, b} = \delta_{ab},
\nonumber\\
h_{j\ell} =
\sigma_{ij}\bar{\sigma}_{k\ell}\,
 z^{i+n}\ol{z}^{k+n},
\nonumber\\
h_{j\, \ell+n} =
\sigma_{ij}\bar{\sigma}_{k\ell}\,
 z^{i+n}\ol{z}^{k},
\nonumber\\
h_{j+n\,\ell} =
\sigma_{ij}\bar{\sigma}_{k\ell}\,
 z^{i}\ol{z}^{k+n},
\nonumber\\
h_{j+n\, \ell+n} =
\sigma_{ij}\bar{\sigma}_{k\ell}\,
 z^{i}\ol{z}^{k},
\label{metriccoeff}
\end{gather}
so that {$(h_{AB})$} and its inverse are of the form
\[
{\bigl(h_{AB}\bigr)}=\begin{pmatrix} h_{ab} &\1_{2n} \\ \1_{2n} & 0\end{pmatrix},
\qquad
{\bigl(h^{AB}\bigr)}=\begin{pmatrix} 0 &\1_{2n} \\ \1_{2n} & -h_{ab}\end{pmatrix}.
\]

Without loss of generality, we can assume that the symmetric matrix {$(\sigma_{ij})$} is diagonalised, as the following lemma shows.
\begin{Lemma}\label{diag:lem}
For any complex symmetric $n\times n$ matrix $\sigma=(\sigma_{ij})$, the pseudo-Riemannian manifold $\bigl(\R^{8n},g\bigr)$ is isometric to a pseudo-Riemannian manifold $\bigl(\R^{8n},\hat g\bigr)$ of the same type defined by a diagonal matrix $(\hat\sigma_{ij})$.
\end{Lemma}

\begin{proof}
Assume that $Q\in \GL(n,\C)$ such that $\hat \sigma=Q^\top \sigma Q$ is diagonal. Let $P=\bigl(\overline{Q}^{-1}\bigr)^\top$. The real coordinate transformation $\phi_Q$ induced by complex linear transformation
\[ \hat{z}^{i}:=Q^i{}_k z^{k},\qquad \hat{z}^{i+n}:=Q^i{}_k z^{k+n},\]
and
\[ \hat{z}^{i+2n}:=P^i{}_k z^{k+2n},\qquad {\hat{z}^{i+3n}}:=P^i{}_k z^{k+3n},\]
maps $\R^{8n}$ to $\R^{8n}$ and pulls back the K\"{a}hler potential $f$ to the K\"{a}hler potential $\hat f= f\circ \phi_Q$,
\[
\hat f\bigl({z}^1,\ldots , {z}^{4n}\bigr)
=\delta_{ab} {z}^a\ol{{z}}^{b+2n}+
\hat\sigma_{ij}\bar{\hat\sigma}_{k\ell}
 \bigl(
 {z}^{i}{z}^{j+n}\ol{{z}}^{k} \ol{{z}}^{\ell+n}
 \bigr).
 \]
 As a consequence, the corresponding pseudo-K\"{a}hler metrics are isometric, $\hat g=\phi_Q^*g$. Note also that the quadratic term $\delta_{ab} {z}^a\ol{{z}}^{b+2n}$ is invariant under $\phi_Q$, i.e.,~$\phi_Q\in \O(4n,4n)$.
\end{proof}

We will now compute the Christoffel symbols for the metric (\ref{metric}) defined by the K\"{a}hler potential~(\ref{potential}). From (\ref{metriccoeff}), the metric is given as
\begin{gather}
g =\delta_{ab} \bigl(\d z^a\d \ol{z}^{b+2n} + \d z^{a+2n}\d \ol{z}^{b}\bigr)
\nonumber\\ \hphantom{g=}{}
 +
\sigma_{ij}\bar{\sigma}_{k\ell}
 \bigl(
 z^{i+n}\ol{z}^{k+n}\d z^{j}\d\ol{z}^{\ell}
 +
 z^{i+n}\ol{z}^{k}\d z^{j}\d\ol{z}^{\ell+n}
 +
 z^{i}\ol{z}^{k+n}\d z^{j+n}\d\ol{z}^{\ell}
 +
 z^{i}\ol{z}^{k}\d z^{j+n}\d\ol{z}^{\ell+n}\bigr)
\nonumber\\ \hphantom{g}{}
 =
\delta_{ab} \bigl(\d z^a\d \ol{z}^{b+2n} + \d z^{a+2n}\d \ol{z}^{b}\bigr)
+
 \sigma_{ij} \d\bigl(z^{i} z^{j+n}\bigr)
 \overline{
 \sigma_{k\ell} \d\bigl(z^{k} z^{\ell+n}\bigr)}
.\label{metriclong}
\end{gather}
In the following, we will write $\partial_A$ for $\frac{\partial}{\partial z^A}$ and $\partial_{\overline{A}}$ for $\frac{\partial}{\partial \overline{z}^A}$.
The {(essential)} Christoffel symbols of a K\"{a}hler metric are given by \cite[Section~IX.5]{ko-no2}
\[\Gamma^C_{AB} =h^{DC}\partial_Ah_{BD}.\]
\big(In fact, \smash{$\Gamma^{\bar{C}}_{\bar{A}\bar{B}}= \overline{\bigl(\Gamma^C_{AB}\bigr)}$} and the mixed Christoffel symbols are zero.\big)
In our situation, this shows~that
\[ \Gamma^A_{a+2n\ B}=0,\]
so that the holomorphic vector fields $\del_{a+2n}$, $a=1,\ldots , 2n$, are parallel on $\bigl(\R^{8n}, g\bigr)$. As~a~consequence, the real vector fields
$\del_{a+2n}+\del_{\overline{a+2n}}$ and $\mathrm{i}{\bigl(\del_{a+2n}-\del_{\overline{a+2n}}\bigr)}$ are parallel as well.

Similarly,
one
checks that
\[\Gamma^{c}_{AB} =0.\]
Furthermore, it is easy to see that
\[
\Gamma^{c+2n}_{i j}= \partial_{i}h_{j\, c}=0, \]
as well as
\[
\Gamma^{c+2n}_{i+n\,j+n}=0.
\]
A similar computation shows that the only non-vanishing Christoffel symbols are
\[
\Gamma^{\ell+2n}_{i+n\, j}
 =
 \partial_{i+n} h_{ j\, \ell}
 =
 \partial_{i+n} \bigl(
 \sigma_{mj}\bar{\sigma}_{k\ell}\,
 z^{m+n}\ol{z}^{k+n}\bigr)
 =
 \sigma_{ij}\bar{\sigma}_{k\ell}\,
\ol{z}^{k+n}
\]
and
 \[ \Gamma^{\ell+3n}_{i+n\, j}
 =
 \partial_{i+n} h_{ j\, \ell+n}
 =
 \partial_{i+n} \bigl(
 \sigma_{mj}\bar{\sigma}_{k\ell}\,
 z^{m+n}\ol{z}^{k} \bigr)
 =
 \sigma_{ij}\bar{\sigma}_{k\ell}\,
\ol{z}^{k} .\]
To summarise, for all $b\in \{1, \ldots, 2n\}$ and $A,B\in \{1, \ldots, 4n\}$, we have
\begin{equation}
\label{LC1}
\nabla_{\partial_A}\partial_{b+2n}=0,\qquad
\nabla_{\partial_A}\partial_b \in \Gamma (\mathrm{span} (\partial_{a+2n})_{a=1,\ldots, 2n}).\end{equation}
The curvature tensor is determined by ${\del_{a+2n}\hook R=0}$ and the following equations:
\begin{alignat*}{3}
&R\bigl(\del_{\overline{i}},\del_{j}\bigr)\del_{k} = 0,\qquad&& R\bigl(\del_{\overline{i}},\del_{j}\bigr)\del_{k+n} = \sigma_{jk}\sum_{\ell}\bar{\sigma}_{i\ell} \del_{\ell+3n},&
\\
&R\bigl(\del_{\overline{i+n}},\del_{j+n}\bigr)\del_{k} = \sigma_{jk}\sum_{\ell}\bar{\sigma}_{i\ell} \del_{\ell+2n},\qquad&& R\bigl(\del_{\overline{i+n}},\del_{j+n}\bigr)\del_{k+n} = 0,&
\\
&R\bigl(\del_{\overline{i}},\del_{j+n}\bigr)\del_{k} = \sigma_{jk}\sum_{\ell}\bar{\sigma}_{i\ell} \del_{\ell+3n},\qquad&& R\bigl(\del_{\overline{i}},\del_{j+n}\bigr)\del_{k+n} = 0,&
\\
&R\bigl(\del_{\overline{i+n}},\del_{j}\bigr)\del_{k} = 0,\qquad&& R\bigl(\del_{\overline{i+n}},\del_{j}\bigr)\del_{k+n} =
\sigma_{jk}\sum_{\ell}\bar{\sigma}_{i\ell} \del_{\ell+2n}.&
\end{alignat*}
Using conjugation and the skew symmetry of $R$, we get
\begin{alignat*}{3}
&R\bigl(\del_{\overline{i}},\del_{j}\bigr)\del_{\overline{k}} = 0,\qquad&& R\bigl(\del_{\overline{i}},\del_{j}\bigr)\del_{\overline{k+n}} = -\bar{\sigma}_{ik}\sum_{\ell}{\sigma}_{j\ell} \del_{\overline{\ell+3n}},&
\\
&R\bigl(\del_{\overline{i+n}},\del_{j+n}\bigr)\del_{\overline{k}} = -\bar{\sigma}_{ik}\sum_{\ell}{\sigma}_{j\ell} \del_{\overline{\ell+2n}}, \qquad&& R\bigl(\del_{\overline{i+n}},\del_{j+n}\bigr)\del_{\overline{k+n}} = 0,&
\\
&R\bigl(\del_{\overline{i}},\del_{j+n}\bigr)\del_{\overline{k}} = 0,\qquad&& R\bigl(\del_{\overline{i}},\del_{j+n}\bigr)\del_{\overline{k+n}} =
-\bar{\sigma}_{ik}\sum_{\ell}{\sigma}_{j\ell} \del_{\overline{\ell+2n}},&
\\
&R\bigl(\del_{\overline{i+n}},\del_{j}\bigr)\del_{\overline{k}} =
-\bar{\sigma}_{ik}\sum_{\ell}{\sigma}_{j\ell} \del_{\overline{\ell+3n}},\qquad&& R\bigl(\del_{\overline{i+n}},\del_{j}\bigr)\del_{\overline{k+n}} =0.&
\end{alignat*}
With this information at hand, we obtain the following proposition, for which we recall
that a~semi-Riemannian manifold is indecomposable if it does not split as a product, not even locally.

\begin{Proposition}\label{kaehlerprop}
The metric $g$ in \eqref{metric} defined by the K\"{a}hler potential \eqref{potential} is locally symmetric, K\"{a}hler of neutral signature, Ricci-flat, but not flat and hence not conformally flat. Moreover, the semi-Riemannian manifold $\bigl(\R^{8n},g\bigr)$ is a symmetric space.
\big(More precisely, $\bigl(\C^4=\R^{8n},g\bigr)$ is a Hermitian symmetric space.\big) If $(\sigma_{ij})$ is non-degenerate, then
$\bigl(\R^{8n},g\bigr)$ is indecomposable.
\end{Proposition}

\bprf
Using the fact~\cite[p.~90]{Moroianu07} that
\[ \mathrm{Ric}(\partial_A,\partial_{\bar B}) = -\partial_A\partial_{\bar B}\log \det (g(\partial_C,\partial_{\bar D}))),\]
we see that since the determinant of $h_{AB}$ is equal to $1$, $g$ is Ricci-flat.
With the above formulas~(\ref{LC1}) for the Levi-Civita connection, we have that
\begin{equation*}
%\label{LC}
\nabla_{\partial_A}\partial_B \in \Gamma (\mathrm{span} (\partial_{a+2n})_{a=1,\ldots, 2n}),\qquad \nabla_{\partial_{\bar{A}}}\partial_{\bar{B}} \in \Gamma \bigl(\mathrm{span} \bigl(\partial_{\overline{a+2n}}\bigr)_{a=1,\ldots, 2n}\bigr).\end{equation*}
Together with the fact that the components of the curvature tensor in the coordinates are constant, this implies that the curvature tensor is parallel, so that $g$ is locally symmetric,
$\nabla R=0$.

In order to show that $\R^{8n}$ with the metric $g$ is a globally symmetric space, we have to show that it is geodesically complete.
The geodesic equations are (we only work with the unbarred components, as the barred ones follow by complex conjugation)
\begin{gather*}
0 = \ddot{\gamma}^{a},\qquad\text{ i.e.,}\quad \gamma^a(t)=p^at+q^a,
\\
0 =
 \ddot{\gamma}^{\ell+2n} {+2}
 \sigma_{ij}\bar{\sigma}_{k\ell}\,
\ol{\gamma}^{k+n} \dot\gamma^{i+n}{\dot\gamma^{j}}
=
\ddot{\gamma}^{\ell+2n} {+2}
 \sigma_{ij}\bar{\sigma}_{k\ell}\,
\bigl(p^{k+n} t+q^{k+n}\bigr)
{p^{i+n}p^{j}},
\\
0 =
 \ddot{\gamma}^{\ell+3n} {+2}
 \sigma_{ij}\bar{\sigma}_{k\ell}\,
\ol{\gamma}^{k} \dot\gamma^{i+n}{\dot\gamma^{j}}
=
 \ddot{\gamma}^{\ell+3n} {+2}
 \sigma_{ij}\bar{\sigma}_{k\ell}\,
\bigl(p^{k} t+q^{k}\bigr){p^{i+n}p^{j}}.
\end{gather*}
Hence the $\gamma^{a+2n}(t) $ are cubic polynomials in $t$ and hence
 defined for all $t\in \R$.
As a consequence, $\bigl(\R^{8n},g\bigr)$ is geodesically complete and a symmetric space of signature $(4n,4n)$.

Now we prove that the holonomy algebra of $\bigl(\R^{8n},g\bigr)$ is indecomposable if $(\sigma_{ij})$ is non-degenerate.
Indecomposability means here that there is no decomposition of the (real) tangent space as a proper sum of two complementary non-degenerate invariant subspaces. This is equivalent to the geometric indecomposabilty formulated in the statement of the proposition.

In virtue of Lemma \ref{diag:lem}, we can assume that $(\sigma_{ij})$ is diagonal
with diagonal elements $\lambda_i\neq 0$.
The holonomy algebra is spanned by the following endomorphisms:
\begin{gather*}
R\bigl(\partial_i + \partial_{\bar i},\partial_j+\partial_{\bar j}\bigr) = R\bigl(\partial_{\bar i}, \partial_j\bigr) - R\bigl(\partial_{\bar j}, \partial_i\bigr),\\
R\bigl(\partial_i + \partial_{\bar i},\mathrm{i}\bigl(\partial_j-\partial_{\bar j}\bigr)\bigr) = \mathrm{i}\bigl(R\bigl(\partial_{\bar i}, \partial_j\bigr) +
R\bigl(\partial_{\bar j}, \partial_i\bigr)\bigr),\\
R\bigl(\partial_{i+n} + \partial_{\overline{i+n}},\partial_{j+n}+\partial_{\overline{j+n}}\bigr) = R\bigl(\partial_{\overline{i+n}}, \partial_{j+n}\bigr) - R\bigl(\partial_{\overline{j+n}}, \partial_{i+n}\bigr),\\
R\bigl(\partial_{i+n} + \partial_{\overline{i+n}},\mathrm{i}\bigl(\partial_{j+n}-\partial_{\overline{j+n}}\bigr)\bigr) = \mathrm{i}\bigl(R\bigl(\partial_{\overline{i+n}}, \partial_{j+n}\bigr) + R\bigl(\partial_{\overline{j+n}}, \partial_{i+n}\bigr)\bigr),\\
R\bigl(\partial_{i} + \partial_{\overline{i}},\partial_{j+n}+\partial_{\overline{j+n}}\bigr) = R\bigl(\partial_{\overline{i}}, \partial_{j+n}\bigr) - R\bigl(\partial_{\overline{j+n}}, \partial_{i}\bigr),\\
R\bigl(\partial_{i} + \partial_{\overline{i}},\mathrm{i}\bigl(\partial_{j+n}-\partial_{\overline{j+n}}\bigr)\bigr) = \mathrm{i}\bigl(R\bigl(\partial_{\overline{i}}, \partial_{j+n}\bigr) + R\bigl(\partial_{\overline{j+n}}, \partial_{i}\bigr)\bigr).
\end{gather*}
Inspection of the above formulas for these endomorphisms shows that the holonomy algebra $\mathfrak{hol}$ at the origin is represented
in the basis $(\partial_A)$ of the holomorphic tangent space ${U\!:=\!T_0^{1,0}\C^{4n}\cong \C^{4n}}$ by the algebra of $4n \times 4n$ matrices of the form
\[ \begin{pmatrix}0&0\\\ast &0\end{pmatrix},\]
where $\ast$ stands for an arbitrary skew-Hermitian $2n \times 2n$ matrix. We claim that the real vector space $U$ does not admit any
non-trivial decomposition into complementary subspaces invariant under $\mathfrak{hol}$. Let $U_1\subset U$ be a proper invariant subspace. It is either contained in
\[ U':= \mathrm{span}\{ \partial_{a+2n}\mid a=1,\ldots, 2n\},\]
which is precisely the kernel of $\mathfrak{hol}$, or has a nontrivial projection
to $U/U'$. Considering a line $L:=\R v$, where $v\in U_1 \setminus U'$, we see that in the latter case
\[ U'=\mathfrak{hol}\cdot L.\]
The reason is that the unitary group acts transitively on the unit sphere. By the holonomy invariance of $U_1$, this implies that $U'\subset U_1$, so that any
invariant subspace of $U$ is either contained in $U'$ or contains $U'$. This implies that the intersection
of any two invariant subspaces of complementary dimensions is non-trivial.
\eprf
